\documentclass[a4paper,12pt]{article}
\usepackage[utf8]{inputenc}
\usepackage[ngerman]{babel}
\usepackage{amsmath}
\usepackage{parskip}

\setlength{\parindent}{0pt}

\title{Modul 294 - Aufgabenblatt 15}
\author{von Lukas Meier}
\date{Unterricht vom 20.04.2026}

\begin{document}

\maketitle

\section{Ziel dieser Abschlusslektion}

Diese Lektion ist bewusst als kurzer Ausblick gedacht. Sie sollen die absoluten Grundlagen von Vue.js sehen, ohne bereits ein komplettes Framework-Projekt aufzubauen.

\section{Was ist Vue?}

Beantworten Sie kurz in eigenen Worten:
\begin{itemize}
    \item Was ist Vue.js?
    \item Welches Problem versucht ein Frontend-Framework zu lösen?
    \item Warum ist es sinnvoll, zuerst Vanilla JavaScript zu lernen?
\end{itemize}

\section{Vue per CDN einbinden}

Erstellen Sie ein kleines Testprojekt mit mindestens diesen Dateien:
\begin{verbatim}
index.html
main.js
\end{verbatim}

Binden Sie Vue im HTML über CDN ein:

\begin{verbatim}
<script src="https://unpkg.com/vue@3/dist/vue.global.js"></script>
<script src="./main.js"></script>
\end{verbatim}

Legen Sie im HTML ein Root-Element an:

\begin{verbatim}
<div id="app"></div>
\end{verbatim}

\section{Die erste Mini-App}

Erstellen Sie in \texttt{main.js} eine kleine Vue-App mit:
\begin{itemize}
    \item einem Titel
    \item einem Textfeld für einen Namen
    \item einer persönlichen Begrüssung
\end{itemize}

Beispielidee:
\begin{itemize}
    \item \texttt{title = "Vue Crash-Course"}
    \item \texttt{name = ""}
\end{itemize}

\section{Direktiven erkennen}

Verwenden Sie in Ihrer Mini-App mindestens:
\begin{itemize}
    \item \texttt{\{\{ \}\}} für Textausgabe
    \item \texttt{v-model} für ein Eingabefeld
    \item \texttt{@click} oder \texttt{v-on:click} für einen Button
\end{itemize}

Beschriften Sie im Code kurz, welche Zeile welche Aufgabe übernimmt.

\section{Liste anzeigen}

Erweitern Sie Ihre App um ein kleines Array, zum Beispiel:
\begin{itemize}
    \item Lieblingssprachen
    \item Tools
    \item Module aus dem Unterricht
\end{itemize}

Zeigen Sie diese Liste mit \texttt{v-for} im Browser an.

\section{Bedingte Anzeige}

Nutzen Sie zusätzlich \texttt{v-if} oder \texttt{v-else}.

Beispiel:
\begin{itemize}
    \item Wenn die Liste leer ist, soll ein Hinweis erscheinen
    \item Sonst wird die Liste angezeigt
\end{itemize}

\section{Vergleich mit Vanilla JS}

Beantworten Sie kurz:
\begin{itemize}
    \item Was wäre an derselben Aufgabe mit reinem JavaScript aufwendiger?
    \item Was wirkt an Vue einfacher oder angenehmer?
\end{itemize}

\section{Ausblick}

Recherchieren oder notieren Sie zu jedem Punkt einen kurzen Satz:
\begin{itemize}
    \item Komponenten
    \item Vue Router
    \item Pinia
    \item Vite
\end{itemize}

Sie müssen diese Werkzeuge heute nicht praktisch einsetzen. Es geht nur darum, die Begriffe schon einmal gesehen zu haben.

\section{Abschluss}

Die Lektion ist abgeschlossen, wenn:
\begin{itemize}
    \item eine kleine Vue-App im Browser läuft
    \item \texttt{v-model}, \texttt{v-for} und \texttt{v-if} sichtbar verwendet werden
    \item Sie den Unterschied zu Vanilla JS grob erklären können
    \item Sie mindestens zwei weiterführende Vue-Themen benennen können
\end{itemize}

\section{Optional}

Falls Zeit bleibt:
\begin{itemize}
    \item Fügen Sie einen Button hinzu, der neue Listeneinträge ergänzt
    \item Erstellen Sie einen einfachen Counter
    \item Überlegen Sie, wie Ihre Notizen-App als Vue-Projekt aussehen könnte
\end{itemize}

\end{document}
